% !TEX root = master.tex
\chapter{Einleitung}
\label{ch:einleitung}

Diese Dokumentation fasst die beiden im Rahmen der Vorlesung \emph{Web-/App-Engineering~II}
(Kurs \DieKursbezeichnung, Sommersemester~2026) von \textbf{Gruppe2} entwickelten Projekte zusammen.
Beide Projekte wurden als Teamarbeit umgesetzt; die Dokumentation ist bewusst als \emph{ein}
gemeinsames Dokument angelegt und in zwei Teile gegliedert.

\textbf{Teil~I} beschreibt das Projekt~A \emph{GuardX} -- eine datenbankgest\"utzte Web-App mit
Session- und Benutzerverwaltung, Admin-Bereich und mehreren sicherheitsorientierten Werkzeugen.
\textbf{Teil~II} beschreibt das Projekt~B \emph{SunShine~Tours} -- eine responsive Web-Seite f\"ur
ein (fiktives) St\"adte-Tour-Unternehmen mit Echtzeit-Kostenberechnung und dateibasierter
Buchungsverwaltung. W\"ahrend Projekt~A den Schwerpunkt auf Sicherheit, Sessions und
Datenbankzugriffe legt, steht bei Projekt~B die clientseitige Interaktivit\"at sowie die
KI-gest\"utzte Erstellung s\"amtlicher Inhalte und Grafiken im Vordergrund.

\section{Zielsetzung der Dokumentation}
Ziel dieser Dokumentation ist es, beide Anwendungen so zu beschreiben, dass sie auf einem
fremden System nachvollziehbar installiert, gestartet und bewertet werden k\"onnen. Daf\"ur werden
je Projekt die Funktionalit\"at, der Aufbau des Frontends, die Verzeichnis- und ggf.
Datenstruktur sowie alle ben\"otigten Zugangsdaten dargestellt. Besonderer Wert wird auf die
Nachvollziehbarkeit gelegt: F\"ur jede geforderte Funktionalit\"at wird in einer
Zuordnungstabelle angegeben, in welchem Skript sie umgesetzt ist (vgl.
Kapitel~\ref{ch:a-funktionen} bzw.~\ref{ch:b-funktionen}).

\section{Aufbau des Dokuments}
Nach dieser Einleitung gliedert sich das Dokument wie folgt. Teil~I (Projekt~A) beginnt mit einem
Referenzdokument (Ziele, Funktionsrahmen, Projektrahmen, Umsetzungsstrategie sowie Aufgaben und
Zust\"andigkeiten), gefolgt vom Product-Backlog, der Beschreibung von Frontend, Verzeichnisstruktur
und Datenbank, der funktionalen Umsetzung sowie einer Installationsanleitung inklusive Zugangsdaten.
Teil~II (Projekt~B) beschreibt analog \"Uberblick und Umfang, die eingesetzten KI-Systeme und
Werkzeuge, die funktionale Umsetzung, die Verzeichnisstruktur und die Installation. Den Abschluss
bildet ein kurzes Fazit.
